\scanpage{359}
prenant les traces on trouve
\[
  2\lambda=2 \tag{45}
\]

--- ce sont les conditions néc. et suff. sur $\lambda$ pour que
\deleted{\unsure{$\lambda$}}
($\varepsilon$ étant unipot. régulier)
$\lambda\varepsilon$ soit encore \inserted{\unsure{str\textsuperscript{t}}}
\unsure{unip.\textsuperscript{t}} (donc néc.\textsuperscript{t} unipotent régulier).
Si $2$ est inv. sur $S$, \inserted{\uline{plus gén.}
si $2$ est section régulière de $\underline{\mathcal O}_S$,}
\deleted{\unsure{est réduit}}
ces conditions impliquent $\lambda=1$, OK, mais
\deleted{\unsure{dans le cas}}\inserted{en dehors} de cette hypothèse (p.ex. en car. $2$
\deleted{\unsure{quitte}} sur un schéma non réduit) il n'en est plus néc. ainsi. Écrivons
\[
  \lambda=1+\nu
\]
les conditions (44) (45) s'écrivent
\[
  2\nu=0,\quad \nu^2=0, \tag{46}
\]
donc pour que sur $S$ l'application (43) soit injection, il f. et s.
que les \inserted{seules} sections de $S$ qui satisfont (46)
sont la section nulle.

En général, désignant par $\mu'_2$ le sous-schéma
en groupes de $\mu_2$ défini par (45)
(c'est en fait un \inserted{schéma en} groupe infinitésimal absolu,
\deleted{\unsure{sur}} \inserted{fini sur $\mathbb Z$, pas plat sur $\mathbb Z$,}
isom. au spectre de
\[
  \mathbb Z[T]/(T^2-1,2T-2)\simeq
  \mathbb Z[S]/(S^2,2\unsure{S}),
\]
isomorphe aussi au schéma des nbs duaux sur $\mathbb Z$,
relatif au module $\mathbb Z/2\mathbb Z$ sur $\mathbb Z$,
\deleted{\unsure{$\mu'_2\simeq$}}
\[
  \mu'_2\simeq\operatorname{Spec}
  \bigl(\mathbb Z\deleted{\unsure{\varepsilon}}\oplus\underbrace{\mathbb Z/2\mathbb Z}_{\text{idéal de carré nul}}\bigr)
  \quad\ldots
\]
\mspage{358}

% Printed archive footer outside the manuscript frame:
% 358
% https://grothendieck.umontpellier.fr
% Université de Montpellier, UMR CNRS 5149, Case courrier 051, Place Eugène Bataillon
% 34095 Montpellier cedex 5 - France

\scanpage{360}
\mspage{481}
Il se trouve que $\mu'_{2S}$ opère
\inserted{librement par translations}
sur le schéma des unipotents réguliers \inserted{\unsure{(ou non)}}
de $G$, et on a
\[
\diagram{\begin{tikzcd}[ampersand replacement=\&]
\underline{\mathrm{Unip.rég}}(G)/\mu'_{2S}
  \arrow[r,"\sim"] \& \underline{\mathrm{Unip.rég}}(PG)
\end{tikzcd}}\ . \tag{47}
\]
Ceci semble à première vue en contradiction avec la \unsure{suite} des Borels
\[
\diagram{\begin{tikzcd}[ampersand replacement=\&,column sep=small,row sep=large,nodes in empty cells]
\deleted{\unsure{1}} \& B \arrow[d,mapsto]
 \& \underline{\mathrm{Borels}}(G)
   \arrow[r,"\sim","B\mapsto B/\mathbb \mathbb G_m"']
   \arrow[d,"\wr"']
 \& \mathrm{Borels}(PG) \arrow[d,"\wr"]
 \& \deleted{\unsure{B}}\widetilde B \arrow[d,mapsto] \\
 \& B_u
 \& \begin{gathered}\text{``Sous-groupes}\\
     \deleted{\unsure{\mathrm U}}\text{unipot. à 1}\\
     \text{param.'' de }G\end{gathered}
     \arrow[r,"\sim"]
 \& \begin{gathered}\text{``Sous-groupes}\\
     \text{unip. à 1 par.''}\\\text{de }PG\end{gathered}
 \& \widetilde B_u \\
 \& \& \text{\unsure{c.à.d.}}\quad L \arrow[r,mapsto,"\sim"] \& p(L)
\end{tikzcd}} \tag{48}
\]
(où $p:G\to PG$ est le morphisme can.\textsuperscript{q})

La difficulté proviendrait-elle du fait que les flèches verticales
(surjectives, par définition de la notion de ``ss-groupes unipot. à 1 paramètre''
de $G$ ou de $PG$) ne seraient pas injectives, en dehors de l'hyp. $2$ inv.
(i.e. car \unsure{rés.}\ $\ne2$)?
J'ai envie de dire que
\[
  \underline B=\operatorname{Norm}(B_u),\quad
  \widetilde B=\operatorname{Norm}(\widetilde B_u) \tag{49}
\]
sont des formules qui permettent de retrouver $B,\widetilde B$ à l'aide
de $B_u,\widetilde B_u$\unsure{,|»} --- et j'ai
\deleted{\unsure{fait}} bien \uline{\unsure{déjà}}
vérifié la première de ces formules, à coups d'alg. lin. dans $M$
(sous la forme : les sous-groupes unipot. à 1 paramètre de $G$
correspondent $1$-$1$ aux sections de $\mathbb P$ sur $S$,
i.e. aux \unsure{ssfibrés} en droites $D$ de $F$\ldots) ---
\mspage{359}

% Printed archive footer outside the manuscript frame:
% 359
% https://grothendieck.umontpellier.fr
% Université de Montpellier, UMR CNRS 5149, Case courrier 051, Place Eugène Bataillon
% 34095 Montpellier cedex 5 - France

\scanpage{361}
mais \unsure{ce même} point dans $PG$, c'est
\deleted{si $u\in\Gamma M^*$ est tel que
$u\widetilde N u^{-1}=\mu N+\lambda$,}
\inserted{pas très trivial.}
\deleted{avec $\mu\in\Gamma\mathbb \mathbb G_m$,}

La difficulté doit plutôt provenir du fait \unsure{qu'une|qu'aucune}
section \unsure{unipotente régulière} $\widetilde\varepsilon$ de $PG$
ne définit pas (ss restriction sur $S$) un sous-groupe unipotent à 1 paramètre
--- alors qu'à une section $\varepsilon$ de
\unsure{unipotents réguliers|unipotente régulière} de $G$
\[
  \varepsilon=1+N,
\]
on définit \inserted{un} par la formule
\[
  \varepsilon(t)=1+tN\ .
\]
On \unsure{notera} que si on remplace $\varepsilon$ par
$\inserted{\varepsilon'=}\lambda\varepsilon$
(avec $\lambda$ satisfaisant (44) (45) $\lambda^2=1$, $2\lambda=2$),
On a
\[
  \varepsilon'=(1+\nu)(1+N)=1+N'
  \qquad N'=(\unsure{1+\nu})N+\nu
\]
et
\[
  \varepsilon'(t)=1+tN'
\]
je dis que si \inserted{$\lambda\ne1$ i.e.}
$\nu\ne0$, \deleted{\unsure{on a alors}} les hom. composés
\[
\diagram{\begin{tikzcd}[ampersand replacement=\&]
\widetilde{\underline\varepsilon},\widetilde{\underline\varepsilon}'\colon
\mathbb G_{\unsure{a|0}} \arrow[r] \& G \arrow[r] \& PG
\end{tikzcd}}
\]
\margin{$\Vert$}
\deleted{\unsure{seront}} associés à $\varepsilon,\varepsilon'$
seront \uline{distincts}, bien qu'ils aient
$\overline{\mathrm m}$ valeur pour $t=1$.
En effet, \deleted{\unsure{on aura}} les dérivées à l'origine
de $\widetilde{\underline\varepsilon},\widetilde{\underline\varepsilon}'$
s'identifient aux composés
\[
\diagram{\begin{tikzcd}[ampersand replacement=\&]
\underline{\mathcal O}_S \arrow[r] \& M \arrow[r] \& M/\underline{\mathcal O}_S\unsure{=|\simeq}P\unsure{\mathfrak g}
\end{tikzcd}}
\]
\mspage{360}

% Printed archive footer outside the manuscript frame:
% 360
% https://grothendieck.umontpellier.fr
% Université de Montpellier, UMR CNRS 5149, Case courrier 051, Place Eugène Bataillon
% 34095 Montpellier cedex 5 - France

\scanpage{362}
\mspage{482}
qui transforment $1$ en $\widetilde N$ resp. en $\widetilde N'$,
(où le $\sim$ désigne ici le passage au quotient \uline{additif}) on a
\[
  \widetilde N'=(1+\nu)\widetilde N
\]
et $\widetilde N$ est une section de $P\unsure{\mathfrak g}$ partout non nulle
--- donc on ne peut avoir \deleted{pas}
\[
  \widetilde N'=\widetilde N\quad\text{i.e.}\quad
  (1+\nu)\widetilde N=\widetilde N\quad\text{i.e.}\quad
  \nu\widetilde N=0
\]
que si $\nu=0$, i.e. $\lambda=1+\nu=1$, cqfd.

En conclusion, il faut considérer que l'on se donne dans $PG$,
non seulement $\widetilde\rho$ et $\widetilde\sigma$
(d'où $\widetilde\varepsilon_0=\widetilde\sigma^{-1}\widetilde\rho$,
$\widetilde\varepsilon_1=\widetilde\rho\widetilde\sigma^{-1}$),
mais aussi des ``sous-groupes à 1 paramètre'' (paramétrés)
\[
\diagram{\begin{tikzcd}[ampersand replacement=\&,nodes in empty cells]
\deleted{\mathbb G_{\unsure{a|0}}}
  \arrow[r,"\deleted{\unsure{\widetilde{\underline\varepsilon}_0}}"] \& \deleted{PG}
\end{tikzcd}}\ ,
\]
\[
\left\{\diagram{\begin{tikzcd}[ampersand replacement=\&,row sep=small]
\widetilde{\underline\varepsilon}_0\colon \mathbb G_{\unsure{a|0}}
  \arrow[r] \& PG \\
\widetilde{\underline\varepsilon}_1\colon \mathbb G_{\unsure{a|0}}
  \arrow[r] \& PG
\end{tikzcd}}\right.\ . \tag{49}
\]
reliés bien sûr par
\[
\left\{\begin{aligned}
\widetilde{\underline\varepsilon}_1
 &=\rho(\widetilde{\underline\varepsilon}_0)
   \bigl(=\sigma(\widetilde{\underline\varepsilon}_0)\bigr)
   \quad\text{i.e.}\\
\widetilde{\underline\varepsilon}_1(t)
 &=\rho\bigl(\widetilde{\underline\varepsilon}_0(t)\bigr)
   =\sigma\bigl(\widetilde{\underline\varepsilon}_0(\unsure{t|u})\bigr)
\end{aligned}\right. \tag{50}
\]
ce qui signifie justement qu'on se \uline{donne}, en plus de
$\widetilde\rho,\widetilde\sigma$, \unsure{un} relèvement\deleted{\illegible{1}}
\inserted{\unsure{unip.\textsuperscript{t}} \inserted{rég.}}
$\varepsilon_0$ de $\widetilde\varepsilon_0$
(\deleted{\unsure{qui|que}} le choix n'est donc pas \unsure{à} priori unique),
et on en déduit $\varepsilon_1$ par
\[
  \deleted{\unsure{\rho}}\quad
  \varepsilon_1=\widetilde\rho(\varepsilon_0)
  =\widetilde\sigma(\varepsilon_0)
\]
\margin{on aura nécessairement une relation
$\sigma^{-1}\rho= -\lambda_0\varepsilon_0$
donc \deleted{$\unsure{\operatorname{int}(\varepsilon_0)|\mathrm{i.e.}\,\sigma(\varepsilon_0)}\,\unsure{\varepsilon_0}$}
$\unsure{\sigma(\varepsilon_0)|\sigma^{-1}(\varepsilon_0)}=\rho(\varepsilon_0)\,\unsure{=}$}
\mspage{361}

% Printed archive footer outside the manuscript frame:
% 361
% https://grothendieck.umontpellier.fr
% Université de Montpellier, UMR CNRS 5149, Case courrier 051, Place Eugène Bataillon
% 34095 Montpellier cedex 5 - France

\scanpage{363}
\mspage{4}
\deleted{\unsure{Et}}
Ainsi, on a ajusté la \deleted{\unsure{précision|position}}
\inserted{description} des données, pour
\deleted{\unsure{ajuste}} pouvoir donner des énoncés de normalisation
précis --- savoir

\uline{Proposition} Soient $\widetilde\rho,\widetilde\sigma$
des sections de $PG$, et $\widetilde{\underline\varepsilon}_0$
un sous-groupe \inserted{additif} \unsure{à} 1-paramètre ``strict''
\margin{(i.e. mono et dont l'image est un $\widetilde B_u$)}
\[
\diagram{\begin{tikzcd}[ampersand replacement=\&]
\widetilde{\underline\varepsilon}_0\colon
\mathbb G_{\unsure{a|0}}
  \arrow[r,phantom,"\unsure{\longrightarrow|\text{---}}" description] \& PG
\end{tikzcd}}\ .
\]
\unsure{tels|tel} que l'on ait \unsure{$\times$|$*$}
\[
  \widetilde\sigma^{-1}\widetilde\rho
   =\widetilde{\underline\varepsilon}_0(1)
   \overset{\mathrm{dfn}}=\widetilde\varepsilon_0 \tag{51}
\]
d'où
\[
\left\{\begin{aligned}
\widetilde\sigma(\widetilde{\underline\varepsilon}_0)
 &=\widetilde\rho(\widetilde{\underline\varepsilon}_0)
   \overset{\mathrm{dfn}}=\widetilde{\underline\varepsilon}_1,
   \quad\text{et}\\
\widetilde\rho\widetilde\sigma^{-1}
 &=\widetilde{\underline\varepsilon}_1(1)
   \overset{\mathrm{dfn}}=\widetilde\varepsilon_1
   =\rho(\widetilde\varepsilon_0)=\sigma(\widetilde\varepsilon_0)
\end{aligned}\right. \tag{52}
\]

Alors :

1) (\unsure{``Pour mémoire''})
$\exists$ uniques $\varepsilon_0,\varepsilon_1$ sections
\inserted{\deleted{\unsure{str\textsuperscript{t}}} unipotentes rég.}
\deleted{telles} de $G$ telles que
(posant $\varepsilon_i=\deleted{\unsure{\mathrm{Id}+}}1+N_i$,
$\underline\varepsilon_i(t)=1+tN_i$) on ait
\[
  \widetilde{\underline\varepsilon}_i(t)=\varepsilon_i(t)^{\sim}\ . \tag{53}
\]

2) \deleted{Supposons que $\varepsilon_0$}
\inserted{On aura alors}
\[
  \unsure{\dot{\varepsilon}_1|\varepsilon_1}=\widetilde\sigma(\varepsilon_0)
    =\widetilde\rho(\varepsilon_0) \tag{54}
\]
\[
\deleted{
  \unsure{\rho}\sigma^{-1}\rho=-\lambda_0\varepsilon_0
  \quad\varepsilon_0=-\lambda_0\sigma^{-1}\rho,\quad
  \varepsilon_1=-\lambda_0\rho\sigma^{-1},
}\LMcancelledtag{\deleted{55}}
\]
\deleted{3) Pour que $\varepsilon_0$ et $\varepsilon_1$ soient
partout indépendants \unsure{en} chaque fibre, i.e. $N$}

3) \unsure{Soient|Étant} $\rho,\sigma$ des relèvements de
$\widetilde\rho,\widetilde\sigma$, on aura alors
\[
  \deleted{\unsure{\rho\sigma^{-1}\rho}}\qquad
  \varepsilon_0=-\lambda_0\sigma^{-1}\rho,\quad
  \varepsilon_1=-\lambda_0\rho\sigma^{-1} \tag{55}
\]
pour $\lambda_0\in\Gamma\mathbb \mathbb G_m$ uniquement déterminé.
Pour un choix convenable de $\rho$ (quand $\sigma$ est déjà choisi)
on peut supposer $\lambda_0=1$
(ce qui détermine $\rho$ de façon unique en
\deleted{\unsure{fonction}} termes de $\sigma$, \unsure{et} inversement)
\mspage{362}

% Printed archive footer outside the manuscript frame:
% 362
% https://grothendieck.umontpellier.fr
% Université de Montpellier, UMR CNRS 5149, Case courrier 051, Place Eugène Bataillon
% 34095 Montpellier cedex 5 - France

\scanpage{364}
\mspage{483}

\unsure{4)} de sorte qu'on a
\[
\tag{56}
\varepsilon_0=-\unsure{\sigma}^{-1}\rho
\qquad(\text{d'où }\varepsilon_1=-\rho\sigma^{-1})
\]

4) $\rho$ et $\sigma$ étant choisis de façon à satisfaire à (56), pour que $\varepsilon_0,\varepsilon_1$ soient \deleted{partout} indépendants \inserted{(en chaque fibre)}, i.e. pour que l'on ait $N_0$ et $N_1$ non homothétiques en \unsure{tte} fibre, il f. et s. qu'on ait
\[
\tag{57}
\operatorname{Tr}(\rho+\sigma)\quad\text{inv.}
\]
(cf. aussi prop. p. 478' pour d'autres conditions équivalentes).

5) Pour que $\varepsilon_0,\varepsilon_1$ soient indépendants \inserted{(en chaque fibre)}, il f. et s. que l'on \deleted{ait} puisse trouver \inserted{(loc. \uline{\unsure{fppf}})} $\rho,\sigma$ au dessus de $\widetilde\rho,\widetilde\sigma$ de façon à satisfaire à (56) \uline{et} à
\[
\tag{57'}
\operatorname{Tr}(\rho+\sigma)=1
\qquad\text{i.e.}\quad
\operatorname{Tr}\rho+\operatorname{Tr}\sigma=1
\]
\margin{$\rho+\sigma$ un projecteur élémentaire $\Longleftrightarrow$}
et cela détermine $\rho,\sigma$ \uline{de façon unique} (en termes de $\widetilde\rho,\widetilde\sigma,\underline{\widetilde\varepsilon}_0$).

\uline{Remarques} 1) Les conditions sur $(\rho,\sigma)$ envisagées dans la prop. p. 478', \unsure{et} la normalisation $\unsure{u|v_1}=1$ i.e. $\operatorname{Tr}(\sigma+\rho)=1$ i.e. $\sigma+\rho$ proj. élém., sont \uline{symétriques} en $(\rho,\sigma)$: elles sont satisfaites pour $(\rho,\sigma)$ ssi elles le sont pour $(\sigma,\rho)$ --- \deleted{\unsure{l'échange de $\rho$ et $\sigma$|l'échange se fait}} \inserted{de façon plus} précise posons
\[
\tag{58}
\rho'=\sigma,\ \sigma'=\rho\qquad
\varepsilon'_0=\underbrace{-\sigma'^{-1}\rho'}_{\unsure{=}-\rho^{-1}\sigma},\quad
\varepsilon'_1=\underbrace{-\rho'\sigma'^{-1}}_{-\sigma\rho^{-1}}
\]
alors on aura
\[
\tag{59}
\varepsilon'_0=\varepsilon_0^{-1},\qquad
\varepsilon'_1=\varepsilon_1^{-1}.
\]
De $\overline{\mathrm m}$, $(\rho^{-1},\sigma^{-1})$ satisfait aux $\overline{\mathrm m}$s conditions \inserted{de la p. 478'}, et \unsure{posant|prenant}

\mspage{363}

% Printed archive footer outside the manuscript frame:
% 363
% https://grothendieck.umontpellier.fr
% Université de Montpellier, UMR CNRS 5149, Case courrier 051, Place Eugène Bataillon
% 34095 Montpellier cedex 5 - France

\scanpage{365}

\[
\tag{60}
\rho''=\rho^{-1},\ \sigma''=\sigma^{-1},\qquad
\varepsilon''_0=\underbrace{-\sigma''^{-1}\rho''}_{-\sigma\rho^{-1}},\quad
\varepsilon''_1=\underbrace{-\rho''\sigma''^{-1}}_{-\rho^{-1}\sigma}
\]
on aura
\[
\tag{61}
\varepsilon''_0=\varepsilon_{\unsure{1|0}}^{-1},\qquad
\varepsilon''_1=\varepsilon_0^{-1},
\]
mais on fera attention qu'on n'aura plus en général $\operatorname{Tr}(\rho''+\unsure{\sigma''})=1$, mais (par un calcul matriciel immédiat)
\[
\tag{62}
\operatorname{Tr}(\rho''+\sigma'')
=\operatorname{Tr}(\rho^{-1}+\sigma^{-1})
=-\frac{1}{\unsure{v|\nu}}
=(\det\rho)^{-1}=(\det\sigma)^{-1}
\]
(en supposant bien sûr \inserted{déjà} $\unsure{v_1|u_1}=1$ i.e. $\operatorname{Tr}(\rho+\sigma)=1$)
--- ce n'est égal à 1 que si
\[
\tag{63}
\det\rho\ (=\det\sigma)=1
\]
--- c'est la condition de la \unsure{section|relation} 3$^\circ$) de $\unsure{(A1)_2|(41)_2}$ (p. 463'), qui signifie donc aussi que la normalisation $\unsure{v_1|u_1}=1$ est stable par la transformation $\rho\mapsto\rho^{-1},\ \sigma\mapsto\sigma^{-1}$.

Par contre, le choix de l'\uline{\unsure{épinglage}} associé à $(\rho,\sigma)$ (\unsure{i.e.} $\widetilde\rho,\widetilde\sigma$) [dans lequel l'accent a été mis sur $\varepsilon_1=-\rho\sigma^{-1}$, qui reçoit la forme simple
\[
\varepsilon_1=\begin{pmatrix}1&0\\1&1\end{pmatrix},
\]
alors que $\varepsilon_0$ \inserted{$=-\sigma^{-1}\rho$} a la forme plus compliquée
\[
\varepsilon_0=\begin{pmatrix}1&\unsure{v_1/v|u_1/v}\\0&1\end{pmatrix}
\qquad\left(=\begin{pmatrix}1&-1\\0&1\end{pmatrix}\text{ dans le cas modulaire}\right)
\]
] fait intervenir implicitement l'ordre dans lequel on se donne $\rho,\sigma$ --- si on le change, on change l'\unsure{épinglage} par le changement de $\varepsilon_{\unsure{1}}$ en $\varepsilon'_1=\varepsilon_1^{-1}$.

En fait, la situation \unsure{offrait} le choix entre \uline{quatre} \unsure{épinglages} ``\unsure{naturels}'', associés à la paire de

\mspage{364}

% Printed archive footer outside the manuscript frame:
% 364
% https://grothendieck.umontpellier.fr
% Université de Montpellier, UMR CNRS 5149, Case courrier 051, Place Eugène Bataillon
% 34095 Montpellier cedex 5 - France

\scanpage{366}
\mspage{484}

sous-groupes à un paramètre $\{L_0,L_1\}$, \deleted{et} \inserted{\unsure{correspondant}} aux quatre sections $\varepsilon_0,\varepsilon_0^{-1},\varepsilon_1,\varepsilon_1^{-1}$ de l'un \deleted{\unsure{ou}} \inserted{et de} l'autre.

On a préféré normaliser à l'aide de $L_1$, à cause de la forme la plus commode de l'équation des lacets (à laquelle j'ai fini par parvenir d'ailleurs après des \unsure{pérégrinations} assez lourdes, où c'est $L_0$ que je m'obstinais à privilégier dans les formules --- cf \S44, la notation définitive n'apparaît qu'en \S45 ---. La préférence donnée à \deleted{$\unsure{\varepsilon}$} $-\rho\sigma^{-1}$ (noté $\varepsilon_1$) à son inverse $-\sigma\rho^{-1}$ (noté $\varepsilon_1^{-1}$) provient de la préférence (justifiée par les faits) donnée à $\rho$ sur $\sigma=\sigma_\infty$, dans la formule des lacets écrite sous la forme
\[
[\alpha,\rho]=[\beta,\sigma]\varepsilon_1(\nu)
\]
au lieu de l'écrire
\[
[\beta,\sigma]=[\alpha,\rho]\varepsilon'_1(\nu)
\qquad(\varepsilon'_1=\varepsilon_1^{-1})
\]
Donc c'est $\varepsilon_1$ qui \unsure{méritait} alors d'avoir la forme matricielle la plus simple possible --- on a pris
\[
\varepsilon_1=\begin{pmatrix}1&0\\1&1\end{pmatrix}\ \ldots
\]

2) En termes des systèmes
\[
\tag{64}
\diagram{\begin{array}{ccc}
\widetilde\ell'_\infty,&\widetilde\ell'_1,&\widetilde\ell'_0\\
\Vert&\Vert&\Vert\\
\widetilde\rho^{-1}&\widetilde\sigma&\widetilde\varepsilon_0
\end{array}}
\quad\text{satisfaisant}\quad
\widetilde\ell'_\infty\cdot\widetilde\ell'_1\cdot\widetilde\ell'_0=1
\]
(i.e. $\widetilde\rho^{-1}\cdot\widetilde\sigma\cdot\widetilde\varepsilon_{\unsure{0}}\unsure{=}1$ i.e. $\widetilde\varepsilon_0=\widetilde\sigma^{-1}\widetilde\rho$), le groupe $\unsure{\mathfrak S'_3|\mathcal G'_3}$ y opère (cf \S42 I) par permutations circulaires, et par changement en $(\widetilde\varepsilon_0^{-1},\widetilde\sigma^{-1},(\widetilde\rho^{-1})^{-1}=\widetilde\rho)$, \unsure{qui} à son tour par permutation circulaire donne $(\widetilde\sigma^{-1},\widetilde\rho,\widetilde\varepsilon_0^{-1})$ ---

\mspage{365}

% Printed archive footer outside the manuscript frame:
% 365
% https://grothendieck.umontpellier.fr
% Université de Montpellier, UMR CNRS 5149, Case courrier 051, Place Eugène Bataillon
% 34095 Montpellier cedex 5 - France

\scanpage{367}

le passage de $(\widetilde\rho^{-1},\widetilde\sigma,\widetilde\varepsilon_0)$ à $(\widetilde\sigma^{-1},\widetilde\rho,\widetilde\varepsilon_0^{-1})$ correspond alors à l'échange $(\rho,\sigma)\mapsto(\sigma,\rho)$. Comment \deleted{\unsure{en appliquer à|en appliquant à} $(\widetilde\rho^{-1},\widetilde\sigma,\widetilde\varepsilon_0)$} interpréter la transformation $(\rho,\sigma)\mapsto(\rho^{-1},\sigma^{-1})$, i.e. $(\widetilde\rho^{-1},\widetilde\sigma,\widetilde\varepsilon_0)\mapsto(\widetilde\rho,\widetilde\sigma^{-1},\widetilde\varepsilon_1)$ ?

Interprétons les triples comme des hom.
\diagram{\begin{tikzcd}[ampersand replacement=\&,column sep=large,row sep=small]
\Pi_{0,3}\arrow[r,"\unsure{\varphi|\psi}"]\&\Gamma PG\\
\ell_0\arrow[r,mapsto]\&\ell'_0\\
\ell_1\arrow[r,mapsto]\&\ell'_1\\
\ell_\infty\arrow[r,mapsto]\&\ell'_\infty
\end{tikzcd}}
on voit que le changement $(\rho,\sigma)\mapsto(\sigma,\rho)$ correspond à la composition de $\unsure{\varphi}$ avec l'autom. de $\Pi_{0,3}$ tel que
\diagram{\begin{tikzcd}[ampersand replacement=\&,column sep=small]
\ell_0\arrow[r,mapsto]\&\ell_1^{-1}\&\ell_1\arrow[r,mapsto]\&\ell_0^{-1}\&\text{donc }\ell_\infty\arrow[r,mapsto]\&\ell_\infty^{-1}
\end{tikzcd}}
i.e. avec \uline{$\tau'_\infty$}. De $\overline{\mathrm m}$ la transformation $(\rho,\sigma)\mapsto(\rho^{-1},\sigma^{-1})$ correspond à la composition avec l'autom.
\diagram{\begin{tikzcd}[ampersand replacement=\&,column sep=small,row sep=small]
\ell_0\arrow[r,mapsto]\&\ell_0^{-1}\&\text{d'où }\ell_\infty\arrow[r,mapsto]\&\ell_0\ell_1\\
\ell_1\arrow[r,mapsto]\&\ell_1^{-1}\&\&{}
\end{tikzcd}},
qui n'est autre que \uline{$\tau_\infty$}$=\sigma_\infty\tau'_\infty$.

Donc le groupe de quatre transformations \deleted{\unsure{qui opère|qu'on a}} sur l'ens. des couples $(\rho,\sigma)$ qu'on a envisagé dans \uline{1$^\circ$)} s'interprète \deleted{\unsure{comme}} \inserted{\deleted{par} comme} le sous-groupe $\simeq\mathbb Z/2\mathbb Z\times\mathbb Z/2\mathbb Z$ engendré par $\tau'_\infty,\tau_\infty$ dans $\operatorname{Aut}_{\unsure{\mathrm{lac}|\mathrm{ext}}}(\Pi_{0,3})=\unsure{\mathfrak S'_{0,3}|\mathcal G_{0,3}}$, qui \deleted{est} stabilise $L_0\subset\Pi_{0,3}$ (donc son action sur \unsure{le|ce} troisième facteur $\widetilde\varepsilon_0=\ell'_{\unsure{0}}$ sera de le conserver ou le remplacer par son inverse, \deleted{\unsure{).}} ce qui aura la vertu de \deleted{\unsure{le|la}} conserver son caractère unipotent, auquel nous tenons !).

\mspage{366}

% Printed archive footer outside the manuscript frame:
% 366
% https://grothendieck.umontpellier.fr
% Université de Montpellier, UMR CNRS 5149, Case courrier 051, Place Eugène Bataillon
% 34095 Montpellier cedex 5 - France

\scanpage{368}
\mspage{485}

\uline{$\overline{\mathrm{VIII}}$}\quad
\uline{Récapitulation sur l'équation}
\[
\tag{1}
[\alpha,\rho]=[\beta,\sigma]\varepsilon_1(\nu).
\]
Les groupes \deleted{$\unsure{\mathcal G_{\rho,\sigma}^{*0}}$}, $S\mathcal G_{\rho,\sigma}^{*}$ et ses variantes.

Le résultat essentiel auquel nous sommes parvenus dans (\uline{$\overline{\mathrm{III}}$}), c'est la possibilité de paramétriser les solutions d'une telle équation \inserted{en $\alpha,\beta,\nu$} (sous des hypothèses convenables sur $\rho,\sigma$) par un groupe $B'_0\times T_\rho\times T_\sigma$, où $B'_0$ est un ``\unsure{pré-Borel|pri-Borel}'' dans $G$ (prenant, pour une base convenable, la forme
\[
B'_0=\left\{\begin{pmatrix}\mu&\lambda\\0&1\end{pmatrix}\;\middle|\;
\begin{array}{l}\mu\in\Gamma \mathbb G_m\\\lambda\in\Gamma \mathbb G_a\end{array}\right.
\]
et $T_\rho,T_\sigma$ sont les \unsure{centralisateurs|stabilisateurs} des sections régulières $\rho,\sigma$ de $G$, par le morphisme
\[
\tag{2}
\left\{\begin{gathered}
\diagram{\begin{tikzcd}[ampersand replacement=\&,column sep=large,row sep=small]
B'_0\times T_\rho\times T_\sigma\arrow[r]\&\mathcal G_{\rho,\sigma}\\
(\underline u,\underline a,\underline b)\arrow[r,mapsto]\&(\underline u\underline a^{-1},\underline u\underline b^{-1},\unsure{\nu|\delta}(\underline u))
\end{tikzcd}}
\end{gathered}\right.
\]
(où le fait que $\nu(\underline u)=\det(\underline u)-1$ n'a pas pour le moment d'importance capitale).

\margin{Nous partons maintenant de $\rho,\sigma,B'_0,\varepsilon_1$ (d'où $L_1$) et nous posons la question à quelles conditions, on peut avoir une paramétrisation (2).}

\deleted{Pour} \inserted{Dire} que ce morphisme \unsure{du premier|des premiers} \unsure{ds} $G\times G\times E'$ \deleted{donne} \inserted{s'envoie} bien dans les solutions de (1), \deleted{il faut et il suffit} \inserted{\unsure{signifie}} aussi qu'on a identiquement
\[
\underline u\rho(\underline u)^{-1}=\underline u\sigma(u)^{-1}\varepsilon_1(\underline\nu(u))
\]
ce qui s'écrit aussi, en simplifiant par $\unsure{\underline u|=}$ et prenant les inverses
\[
\tag{3}
\rho(\underline u)=\varepsilon_1(-\underline\nu(u))\sigma(\underline u).
\]
La possibilité de trouver $\underline\nu:B'_0\longrightarrow \mathbb G_a$ (morphisme de schémas) \unsure{donnera|donnerait}, via (2), un morphisme $B'_0\times T_\rho\times T_\sigma\longrightarrow\mathcal G_{\rho,\sigma}\overset{\mathrm{dfn}}{=}$ schéma

\mspage{367}

% Printed archive footer outside the manuscript frame:
% 367
% https://grothendieck.umontpellier.fr
% Université de Montpellier, UMR CNRS 5149, Case courrier 051, Place Eugène Bataillon
% 34095 Montpellier cedex 5 - France

\scanpage{369}

des solutions de (1), équivaut donc à celle
\deleted{\unsure{que}} \deleted{pour tout} $\sigma$ et $\rho$
\deleted{\unsure{aient}} $\overline{\mathrm m}$ action
sur $B'_0$, \uline{modulo} multiplication à \unsure{g.} par
\deleted{$\unsure{L_1}$} le sous-groupe \deleted{$\unsure{L_1}$}
\inserted{\unsure{unip\textsuperscript{t}}} à un paramètre $L_1$ qui
intervient dans l'écriture même de l'équation (1) que nous étudions.
Celà se conçoit difficilement si on n'admet que
$L_1\subset\rho(B'_0)\overset{\mathrm{df}}=B'_1$, ce qui implique d'ailleurs
$\sigma(B'_0)\subset B'_1$ (et par un argument connu $\sigma(B'_0)=B'_1$)
et $L_1=(B'_1)_u$. Inversement, si on a
$\rho(B'_0)=\sigma(B'_0)$, \deleted{\unsure{alors pour tout}}
\inserted{i.e. $\sigma^{-1}$}\deleted{$\unsure{\rho(-)}$}$\rho$ normalise $B'_0$,
\deleted{alors} i.e. $\sigma^{-1}\rho$ section de $B_0$, alors
$\sigma(u)=\sigma u\sigma^{-1}$ et $\rho(u)=\rho u\rho^{-1}$ sont deux sections
de $B'_1$ qui ont $\overline{\mathrm m}$ déterminant -- \unsure{donc ablt\textsuperscript{n}},
donc qui coïncident bien modulo $(B'_1)_u$.

Ainsi on voit que la propriété essentielle pour obtenir des solutions de (1) via (2),
c'est d'avoir
\[
(4)\qquad (\sigma^{-1}\rho)(B'_0)=B'_0\qquad\text{i.e.}\quad
\sigma^{-1}\rho\in\Gamma B_0\, ..
\]
\margin{\uline{NB} Cette condition \unsure{(si} $\sigma^{-1}\rho$ régulier,
\unsure{i.e.} $\rho$ et $\sigma$ \unsure{ne sont homoth.}
\unsure{en aucune fibre)}, détermine pratiquement $B_0$
-- \deleted{\unsure{détermine}} \unsure{sur un corps, il y a une}
\uline{\unsure{ambiguïté}} \unsure{possible} \unsure{entre|d'ordre} \uline{deux}
(les \uline{\unsure{deux}} Borels contenant \unsure{l'élément en q.}).}

Pour l'unicité de la représentation d'un triple $(\alpha,\beta,\nu)$ sous la forme
$(\underline u\underline a^{-1},\underline u\underline b^{-1},\underline\nu(\underline u))$,
il faut \deleted{supposer} \inserted{et il} s. qu'on ait
\[
(5)\qquad T_\sigma\cap T_\rho\cap B'_0=\{1\},
\]
\margin{Il suffit p. ex qu'on ait $T_\rho\cap T_\sigma=\unsure{\mathbb \mathbb G_m}$
et $\mathbb \mathbb G_m\cap B'_0=\{1\}$.}
\deleted{ce qui équivaut aussi à \unsure{qu'on ait}
\[
(6)\qquad T_\sigma\cap\unsure{\rho} B'_0\unsure{=}
\]}
l'indétermination est

\mspage{368}

% Printed archive footer outside the manuscript frame:
% 368
% https://grothendieck.umontpellier.fr
% Université de Montpellier, UMR CNRS 5149, Case courrier 051, Place Eugène Bataillon
% 34095 Montpellier cedex 5 - France

\scanpage{370}
\mspage{486}

dans la forme
\diagram{\begin{tikzcd}[ampersand replacement=\&]
(\underline u,\underline a,\underline b) \arrow[r,mapsto] \&
(\underline u l,\underline a l,\underline b l)
\end{tikzcd}},
avec $l$ une section de $T_\rho\cap T_\sigma\cap B'_0$.
Reste la question de la surjectivité de (2) -- toute solution de (1)
s'écrit-elle bien sous la forme
$(u a^{-1},\underline u b^{-1},\nu(\underline u))$ ?
Il revient au $\overline{\mathrm m}$ (si $\rho,\sigma$ sont réguliers)
de dire que \deleted{\unsure{il}} l'on a surjectivité dans
\[
(6)\qquad
\diagram{\begin{tikzcd}[ampersand replacement=\&]
B'_0 \arrow[r] \& \mathcal X_{\rho,\sigma} \\
\underline u \arrow[r,mapsto] \&
([\underline u,\rho],[\underline u,\sigma],\underline\nu(\underline u))
\end{tikzcd}}
\]
où $\mathcal X_{\rho,\sigma}$ est le schéma des solutions simultanées
\deleted{en} en \deleted{$\unsure{\underline\nu,}$}
$\zeta,\eta,\nu$ de
\[
(7)\qquad
\det\zeta=\det\eta=1,\quad
\operatorname{Tr}\zeta\rho=\operatorname{Tr}\rho,\quad
\operatorname{Tr}\eta\sigma=\operatorname{Tr}\sigma,\quad
\zeta=\eta\varepsilon_1(\nu).
\]

C'est essentiellement pour obtenir cette surjectivité que nous avions été amenés
dans la section (II) à faire les hypothèses draconiennes
(p.\ 435',436) qui s'expriment, au prix de un étale,
par les conditions \deleted{\unsure{que|pour}}
\[
(8)\qquad
\left\{\begin{array}{ll}
\mathrm{a)} & \sigma^{-1}\rho\text{ section de }(B_0)_u=L_0
\quad(\text{et pas seulement de }B_0)\\
\mathrm{b)} & B_0\text{ et }B_1=\rho(B_0)=
\deleted{\unsure{\rho}}\inserted{\sigma}(B_0)
\text{ en ``position générale'',}\\
&\text{i.e. }L_0\text{ et }L_1=\rho(L_0)=\sigma(L_0)
\text{ indépendants;}\\
&\text{(en chaque fibre)}
\end{array}\right.
\]
mais il n'est pas clair si ces conditions
(suggérées par des calculs antérieurs, et donnant lieu en effet à
\unsure{des|de} très jolis calculs) sont vraiment indispensables pour le résultat de paramétrisati-

\mspage{369}

% Printed archive footer outside the manuscript frame:
% 369
% https://grothendieck.umontpellier.fr
% Université de Montpellier, UMR CNRS 5149, Case courrier 051, Place Eugène Bataillon
% 34095 Montpellier cedex 5 - France

\scanpage{371}

Dans notre contexte il est \inserted{\unsure{assez}} naturel de supposer
$\varepsilon_0=\unsure{\sigma^{-1}\rho|-\sigma^{-1}\rho}$ unipotent
(ou du moins \unsure{quasi-unipotent}, i.e. \deleted{\unsure{un par}}
un itéré $\varepsilon_0^n$ unipotent) mais la condition b) de (8), plus subtile,
apparaît comme moins impérative. La suite montrera si cette restriction est judicieuse.

Dans le cas où on considère une équation
\[
(9)\qquad
\underbrace{[\alpha,\rho]}_{\zeta}=
\varepsilon\underbrace{[\beta,\sigma]}_{\unsure{\eta}}
\varepsilon_1(\nu)
\qquad\deleted{\unsure{\varepsilon}}
\]
avec un paramètre $\varepsilon$ (pouvant notamment prendre la valeur importante $-1$)
$\in\Gamma\mu_2$, \unsure{on|où} nous avons pour la première fois dans ces calculs
détruit la symétrie entre $\rho$ et $\sigma$, en introduisant la condition
$\operatorname{Tr}\sigma=0$ comme condition de stabilité du schéma
\inserted{$\mathcal X_\sigma$} des $\eta$ satisfaisant
\[
\det\eta=1,\qquad\operatorname{Tr}\eta\sigma=\operatorname{Tr}\sigma
\]
par $\mu_2$. Le résultat essentiel auquel nous tenions, était l'isomorphisme
\[
(10)\qquad
\diagram{\begin{tikzcd}[ampersand replacement=\&,column sep=large]
B'_0\times T_\rho\times N'_\sigma \arrow[r,"\sim"] \& \mathcal G_{\rho,\sigma}^{*}\\
(\underline u,\underline a,\underline b) \arrow[r,mapsto] \&
(\alpha=\underline u\underline a^{-1},\beta=\underline u\underline b^{-1},
\mu=\det\underline u,\varepsilon=\varepsilon(\underline b))
\end{tikzcd}}
\]
où $N'_\sigma\subset N_\sigma$ est maintenant le normalisateur du $T_\sigma$
-- défini par la condition
\[
(11).\qquad \underline b\,\sigma(\underline b)^{-1}=
\varepsilon(\underline b)\deleted{\unsure{\sigma}}
\qquad\text{i.e.}\quad\underline b(\sigma)=\varepsilon(\underline b)\sigma
\]

\mspage{370}

% Printed archive footer outside the manuscript frame:
% 370
% https://grothendieck.umontpellier.fr
% Université de Montpellier, UMR CNRS 5149, Case courrier 051, Place Eugène Bataillon
% 34095 Montpellier cedex 5 - France

\scanpage{372}
\mspage{487}

Que le morphisme (10) s'envoie bien dans $\mathcal G_{\rho,\sigma}^{*}$,
est trivial par définition de $N'_\sigma$ et via le fait analogue pour (2)
-- sans hypothèse du genre $\operatorname{Tr}\sigma=0$.
Mais la possibilité d'obtenir (\unsure{loc\textsuperscript{t}})
n'importe quel $\varepsilon$ via un $\underline b\in\Gamma N'_\sigma$
-- ou simplement $\varepsilon=-1$, si 2 inv.) exige justement
$\operatorname{Tr}\sigma=0$.

Bien entendu, nous allons considérer $\mathcal G_{\rho,\sigma}^{*}$ comme un schéma en groupes,
par transport de \unsure{structure} de la structure de groupe produit de
$B'_0\times T_\rho\times N'_\sigma$.
On y distingue le sous-groupe
\[
(12)\qquad \mathcal G_{\rho,\sigma}^{\unsure{\flat|b}*}(0)
\quad\text{correspondant à :}\quad T'_0\times T_\rho\times N'_\sigma
\]
où
\[
(13)\qquad T_0=B'_0\cap B'_1
\]
est le tore maximal privilégié de $B'_0$, qui
\deleted{pour notre} l'épinglage canonique se décrit comme
\[
(14)\qquad T_0=\left\{\left(\begin{array}{cc}\mu&0\\0&1\end{array}\right.\unsure{\big\}|\big)}
\,\middle|\,\mu\in\Gamma\mathbb \mathbb G_m\right\}.
\]
On pose de $\overline{\mathrm m}$
\[
(15)\qquad \mathcal G_{\rho,\sigma}^{\unsure{\flat|b}}(0)
=\mathcal G_{\rho,\sigma}\cap\mathcal G_{\rho,\sigma}^{*\unsure{\flat|b}}(0),
\quad\text{correspondant à : }T'_0\times T_\rho\times T_\sigma
\]

\mspage{371}

% Printed archive footer outside the manuscript frame:
% 371
% https://grothendieck.umontpellier.fr
% Université de Montpellier, UMR CNRS 5149, Case courrier 051, Place Eugène Bataillon
% 34095 Montpellier cedex 5 - France

\scanpage{373}

On aura donc des structures de produit semi-direct
\[
(16)\qquad
\mathcal G_{\rho,\sigma}\simeq
\mathcal G_{\rho,\sigma}^{\unsure{\flat|b}}(0)\cdot\Lambda_0,
\qquad
\mathcal G_{\rho,\sigma}^{*}\simeq
\mathcal G_{\rho,\sigma}^{*\unsure{\flat|b}}(0)\cdot\Lambda_0,
\]
où
\[
(17)\qquad
\left\{\begin{array}{l}
\displaystyle\Lambda_0\subset\mathcal G_{\rho,\sigma},\quad
\Lambda_0=\left\{\left(
\alpha=\underset{\substack{\Vert\\\varepsilon_0^{\unsure{n|u}}}}{\lambda},
\beta=\underset{\substack{\Vert\\\varepsilon_0^{\unsure{n|u}}}}{\lambda},
\mu=1\right)\,\middle|\,\lambda\in\Gamma L_0\right\}\\
\hfill(\unsure{n|u}\in\Gamma\mathbb \mathbb G_a)\\[4pt]
\text{correspondant à :}\quad L_0\times\{1\}\times\{1\}
\subset B'_0\times T_\rho\times T_\sigma.
\end{array}\right.
\]

Nous allons enfin considérer un autre sous-\inserted{schéma en} groupes de
$\mathcal G_{\rho,\sigma}^{*}$, par les conditions
\[
(18)\qquad
\left\{\begin{array}{l}
\det\alpha\in\unsure{\Gamma\mu_2}\quad\text{i.e. }(\det\alpha)^2=1\\
\det\beta=1
\end{array}\right.
\qquad\left|\quad\text{d'où }S\mathcal G_{\rho,\sigma}^{\unsure{*|\flat}}
\subset\mathcal G_{\rho,\sigma}^{*}\right.
\]

\deleted{Soit donc} \inserted{\unsure{Rem.}}
\deleted{$\unsure{\varepsilon'\longleftarrow(\alpha,\beta),|\varepsilon'(\alpha,\beta),}$}
$S\mathcal G_{\rho,\sigma}^{*}$ peut s'interpréter comme le schéma des systèmes
\[
(\alpha,\beta,\mu,\varepsilon,\varepsilon')
\]
satisfaisant
\[
(19)\qquad [\alpha,\rho]=\varepsilon[\beta,\sigma]\varepsilon_1(\mu-1),
\quad\det\alpha=\varepsilon',\quad
\inserted{\det\beta=1,}\quad
\underline{\varepsilon^2=\varepsilon'^2=1}
\]
\margin{NB $\varepsilon^2=1$ résulte déjà de la première relation.}

En termes des systèmes
\[
(\underline u,\underline a,\underline b)\in\Gamma B'_0\times T_\rho\times N'_\sigma
\]
les équations (19) (dont le \unsure{premier|première} est automatique), se réduisent à :
\[
(20)\qquad\det\underline u=\det\underline b=\varepsilon'\det\underline a,
\qquad\varepsilon'^2=1,
\]

\mspage{372}

% Printed archive footer outside the manuscript frame:
% 372
% https://grothendieck.umontpellier.fr
% Université de Montpellier, UMR CNRS 5149, Case courrier 051, Place Eugène Bataillon
% 34095 Montpellier cedex 5 - France

\scanpage{374}
\mspage{488}
qui \unsure{montrent} que $S\mathcal{G}_{\rho,\sigma}^{*}$ est \inserted{bien} un \uline{sous-groupe}
de $\mathcal{G}_{\rho,\sigma}^{*}$. Ce sous-groupe contient $\Lambda_0$
(\unsure{correspondant} : $L_0\times\unsure{\{1\}\times\{1\}|\{(1,1)\}}\subset B'_0\times T_\rho\times N'_\sigma$) --- qui
correspond au cas où \deleted{\unsure{l'on a}}
\[
\det\underline{u}=\det\underline{b}=\det\underline{a}=1\qquad(\varepsilon'=1),
\tag{21}
\]
et on a une structure de produit $\unsure{\frac12}$ direct
\[
S\mathcal{G}_{\rho,\sigma}^{*}\simeq S\mathcal{G}_{\rho,\sigma}^{*\unsure{\flat|b}}(0)\mathbin{.}\Lambda_0
\tag{22}
\]
où
\[
S\mathcal{G}_{\rho,\sigma}^{*}(0)=S\mathcal{G}_{\rho,\sigma}^{*}\cap\mathcal{G}_{\rho,\sigma}^{*\unsure{\flat|b}}(0)
\tag{23}
\]
correspond au \unsure{ssgroupe|s/groupe} \deleted{de} $\deleted{\unsure{B'_0}}\unsure{T_0|T'_0}\times T_\rho\times N'_\sigma$
\deleted{de $B'_0\times T_\rho\times N_\sigma$} défini par les $\overline{\mathrm m}$s équations
(20). Bien sûr, on aura aussi
\[
S\mathcal{G}_{\rho,\sigma},\quad S\mathcal{G}_{\rho,\sigma}^{\unsure{\flat|b}}(0)\quad(\text{sans }*)
\tag{24}
\]
\deleted{comme} \inserted{intersections} de $S\mathcal{G}_{\rho,\sigma}^{*}$, $S\mathcal{G}_{\rho,\sigma}^{\unsure{\flat|b}*}(0)$ avec
$S\mathcal{G}_{\rho\sigma}$ --- déduit de l'expression
de ceux-ci en termes de $(B'_0,\unsure{T_0|T'_0}),T_\rho,N'_\sigma$
en remplaçant $N'_\sigma$ par $T_\sigma$.

Pour expliciter la structure des $S\mathcal{G}_{\rho,\sigma}^{*}$,
on est ramené via (22) \uline{à étudier le}
groupe
\fbox{\diagram{\begin{tikzcd}[ampersand replacement=\&]
(25)\quad S\mathcal{G}_{\rho,\sigma}^{*\unsure{\flat|b}}(0) \arrow[r, hook] \& \unsure{T_0|T'_0}\times T_\rho\times N'_\sigma
\end{tikzcd}}}
\mspage{373}

% Printed archive footer outside the manuscript frame:
% 373
% https://grothendieck.umontpellier.fr
% Université de Montpellier, UMR CNRS 5149, Case courrier 051, Place Eugène Bataillon
% 34095 Montpellier cedex 5 - France

\scanpage{375}
décrit par les équations (20). Mais on a
\diagram{\begin{tikzcd}[ampersand replacement=\&]
T_0 \arrow[r,"\sim"] \& \mathbb{G}_m \\
\underline{u} \arrow[r,mapsto] \& \det\underline{u}
\end{tikzcd}},
ce qui implique que la projection
\diagram{\begin{tikzcd}[ampersand replacement=\&,column sep=large]
S\mathcal{G}_{\rho,\sigma}^{*\unsure{\flat|b}}(0) \arrow[r] \& T_\rho\times N'_\sigma \\
(\underline{u}\underline{a}^{-1},\underline{u}\underline{b}^{-1},\det\underline{u},\varepsilon=\varepsilon(\underline{b})) \arrow[r] \arrow[d,phantom,"\wr" description] \& (\underline{a},\underline{b}) \\
(\underline{u},\underline{a},\underline{b}) \&
\end{tikzcd}}
est mono, et induit un iso. de
$S\mathcal{G}_{\rho,\sigma}^{*}(0)$ sur le \unsure{ssgroupe|s/groupe} de $T_\rho\times N'_\sigma$
\unsure{défini|définis} par l'équation
\[
\det\underline{b}=\varepsilon'\det\underline{a},\qquad \varepsilon'{}^2=1
\tag{21}
\]
on aura
\[
(\det\underline{b})^2=(\det\underline{a})^2.
\]

En termes des homomorphismes (\unsure{déterminants}, en somme)
\diagram{\begin{tikzcd}[ampersand replacement=\&]
T_\rho \arrow[r,"\delta_\rho"] \& \mathbb{G}_m, \& N'_\sigma \arrow[r,"\unsure{\delta_\sigma}"] \& \mathbb{G}_m
\end{tikzcd}}
cela revient à prendre l'image inverse de la
diagonale \inserted{de $\mathbb{G}_m\times \mathbb{G}_m$} par $\delta_\rho^2,\delta_\sigma^2$.

On voit tout de suite (comme $\delta_\rho,\delta_\sigma$
sont épi) que l'hom
\diagram{\begin{tikzcd}[ampersand replacement=\&,column sep=large]
(22)\quad S\mathcal{G}_{\rho,\sigma}^{*\unsure{\flat|b}}(0) \arrow[r] \& \mu_2\times\mu_2 \\
(\alpha,\beta,\mu,\varepsilon,\varepsilon') \arrow[r,mapsto] \& (\varepsilon,\varepsilon') \\
\text{i.e.}\quad(\underline{u},\underline{a},\underline{b}) \arrow[r,mapsto] \& \left(\underline{\varepsilon}(\underline{b}),\underbrace{\det\underline{b}/\det\underline{a}}_{\det(\underline{b}\underline{a}^{-1})}\right)
\end{tikzcd}}
est épi, et \deleted{le noyau} donne donc
naissance à une suite exacte
\mspage{374}

% Printed archive footer outside the manuscript frame:
% 374
% https://grothendieck.umontpellier.fr
% Université de Montpellier, UMR CNRS 5149, Case courrier 051, Place Eugène Bataillon
% 34095 Montpellier cedex 5 - France

\scanpage{376}
\mspage{489}
\diagram{\begin{tikzcd}[ampersand replacement=\&,column sep=small]
(23)\quad1 \arrow[r] \& S\mathcal{G}_{\rho,\sigma}^{*\circ\unsure{\flat|b}}(0) \arrow[r] \& S\mathcal{G}_{\rho,\sigma}^{*\unsure{\flat|b}}(0) \arrow[r] \& \mu_2\times\mu_2 \arrow[r] \& 1
\end{tikzcd}}
où $S\mathcal{G}_{\rho,\sigma}^{\unsure{\circ\flat|\circ}}(0)$ est isom. au sous-schéma de
$T_\rho\times T_\sigma$ défini par l'équation
\[
\det\underline{a}=\det\underline{b}
\tag{24}
\]
et se dévisse donc à son tour en
\diagram{\begin{tikzcd}[ampersand replacement=\&,column sep=small]
(25)\quad1 \arrow[r] \& ST_\rho^{\unsure{\flat|b}}\unsure{\times} ST_\sigma^{\unsure{\flat|b}} \arrow[r] \& S\mathcal{G}_{\rho,\sigma}^{\circ\unsure{\flat|b}}(0) \arrow[r] \& \mathbb{G}_m \arrow[r] \& 1 \\
\& \& \deleted{\unsure{(\underline{u},|\underline{u}}}(\underline{a},\underline{b}) \arrow[r,mapsto] \& \deleted{\det\underline{u}}\det\underline{a}=\det \&
\end{tikzcd}}
où on a posé
\[
\left\{
\begin{aligned}
ST_\rho^{\unsure{\flat|b}}&=\operatorname{Ker}\left(T_\rho\xrightarrow{\delta_\rho}\mathbb{G}_m\right)\\
ST_\sigma^{\unsure{\flat|b}}&=\operatorname{Ker}\left(T_\sigma\xrightarrow{\delta_\sigma}\mathbb{G}_m\right)
\end{aligned}
\right.
\tag{26}
\]

Si par exemple $\rho$ \inserted{(\unsure{non-s|non-g})} est $\frac12$ simple, donc $T_\rho$ un
tore maximal, alors $ST_\rho^{\unsure{\flat|b}}$ est lui-$\overline{\mathrm m}$ un tore \inserted{(de \unsure{dim 1})} ---
\unsure{itou} pour $\sigma$. Ainsi, \deleted{\illegible{1}} dans le ``cas modulaire''
\[
\rho=\begin{pmatrix}0&1\\-1&1\end{pmatrix},\qquad
\sigma=\begin{pmatrix}0&-1\\1&0\end{pmatrix},
\]
$ST_\rho^{\unsure{\flat|b}}$ est un
tore \deleted{de} sauf en car. 3, et $ST_\sigma^{\unsure{\flat|b}}$ un tore sauf en
car. 2. \uline{Donc}, \uline{sauf en car. \inserted{2,3}}, $S\mathcal{G}_{\rho,\sigma}^{\circ\unsure{\flat|b}}(0)$ est un
\uline{tore de dim 3}.

Pour terminer de déterminer la structure
de $S\mathcal{G}_{\rho,\sigma}^{*}$ en termes de $\mathcal{G}_{\rho,\sigma}^{*\unsure{\circ}}(0)$ et
de la décomposition (22) en produit semi-direct, il faut \deleted{préciser}\inserted{expliciter} l'opération de
\mspage{375}

% Printed archive footer outside the manuscript frame:
% 375
% https://grothendieck.umontpellier.fr
% Université de Montpellier, UMR CNRS 5149, Case courrier 051, Place Eugène Bataillon
% 34095 Montpellier cedex 5 - France

\scanpage{377}
$S\mathcal{G}_{\rho,\sigma}^{*\circ}(0)$ sur $\Lambda_0$, en \unsure{interprétant}
$S\mathcal{G}_{\rho,\sigma}^{*\circ}(0)$ comme sous-groupe de $T_\rho\times N'_\sigma$,
et $\Lambda_0$ comme \inserted{$L_0$} $\simeq \mathbb{G}_{\unsure{a|0}}$ par $\lambda\mapsto\varepsilon_0(\lambda)$.

Or les relations
\[
\deleted{\underline{u}\unsure{\varepsilon_0(\lambda)}\underline{u}^{-1}=}
\]
\[
\underline{u}(\varepsilon_0(\lambda))\underline{u}^{-1}
 =\varepsilon_0(\det(\underline{u})\lambda),\qquad
\det(\underline{u})=\det(\underline{b})=\delta_\sigma(\underline{b})
\]
montrent que l'opération se fait via
le caractère composé
\diagram{\begin{tikzcd}[ampersand replacement=\&]
(27)\quad S\mathcal{G}_{\rho,\sigma}^{*\unsure{\circ|\flat}}(0) \arrow[r] \& N'_\sigma \arrow[r,"\delta_\sigma"] \& \mathbb{G}_m
\end{tikzcd}}, OK.

Pour la suite, nous aurons besoin
d'un hom. canonique
\diagram{\begin{tikzcd}[ampersand replacement=\&,column sep=large,row sep=large]
(28)\quad\mathcal{G}_{\rho,\sigma}^{*} \arrow[rr,"\unsure{\theta|\Theta}"] \arrow[dr,shift right=0.5ex,"{\deleted{\begin{array}{cccc}(\unsure{\alpha},&\beta,&\mu,&\unsure{\varepsilon})\\\unsure{\Vert|=}&\unsure{\Vert|=}&&\\\underline{u}\underline{a}^{-1}&\underline{u}\underline{b}^{-1}&&\end{array}}}"'] \& \& G \\
\& B'_0\times T_\rho\times\unsure{N'_\sigma|T'_\sigma} \arrow[ul,shift right=0.5ex] \arrow[ur,dashed] \&
\end{tikzcd}}
qui s'identifie à la \uline{première} \uline{projection}
\diagram{\begin{tikzcd}[ampersand replacement=\&]
(\underline{u},\underline{a},\underline{b}) \arrow[r,mapsto] \& \underline{u}
\end{tikzcd}}
du produit, mais \unsure{considérée|considéré} comme
un morphisme dans $G$ grâce à l'inclusion
\diagram{\begin{tikzcd}[ampersand replacement=\&]
B'_0 \arrow[r,phantom,"{\unsure{\hookrightarrow|\subset}}" description] \& G.
\end{tikzcd}}
On posera aussi \unsure{souvent}
\unsure{(}
\[
\underline{u}=\unsure{\theta|\Theta}(\alpha,\beta,\mu,\varepsilon),
\]
\unsure{donc|d'où} \deleted{si $\det$} on aura
\mspage{376}

% Printed archive footer outside the manuscript frame:
% 376
% https://grothendieck.umontpellier.fr
% Université de Montpellier, UMR CNRS 5149, Case courrier 051, Place Eugène Bataillon
% 34095 Montpellier cedex 5 - France
